\setauthor{Sebastian Radić \& Luis Schörgendorfer}

\section{Ausgangssituation und Motivation}
Die PROGRAMMIERFABRIK GmbH verarbeitet im Tagesgeschäft regelmäßig eingehende Rechnungen.
Bislang geschah das manuell: Rechnungsdaten wurden von Hand in das interne System eingetippt,
was zeitaufwendig und fehleranfällig war. Mit wachsendem Rechnungsvolumen wurde klar,
dass dieser Prozess langfristig nicht skalierbar ist.

Gleichzeitig ist die elektronische Rechnungsverarbeitung kein rein betriebsinternes Thema mehr.
Die europäische Richtlinie\footnote{Vgl. Europäisches Parlament und Rat der EU: \textit{Richtlinie 2014/55/EU über die elektronische Rechnungsstellung bei öffentlichen Aufträgen}, \url{https://eur-lex.europa.eu/legal-content/DE/TXT/?uri=CELEX:32014L0055}, letzter Zugriff am 20.01.2026}
2014/55/EU verpflichtet öffentliche Auftraggeber zur Annahme elektronischer Rechnungen,
und in Österreich gilt seit 2014 die verpflichtende elektronische Rechnungslegung gegenüber
Bundesbehörden über das Unternehmensserviceportal (USP) im Format ebInterface.\footnote{Vgl. AUSTRIAPRO: \textit{ebInterface -- Der österreichische Standard für die elektronische Rechnung}, \url{https://www.ebinterface.at/}, letzter Zugriff am 20.01.2026}
Diese regulatorischen Vorgaben verstärken den Druck auf Unternehmen, ihre Rechnungsverarbeitung
zu automatisieren.

Der \textit{SmartBillConverter} ist die Antwort auf genau diese Kombination aus betrieblichem
Aufwand und gesetzlichem Druck. Das System nimmt eine eingehende Rechnung als PDF oder Bild
entgegen, extrahiert den Inhalt automatisch und liefert am Ende ein valides XML im Format
ebInterface~6.1 oder ZUGFeRD~2.3 zurück.

\section{Problemstellung und Rahmenbedingungen}
Das zentrale Problem ist, dass Rechnungen, die von Lieferanten einlangen, kein einheitliches
Format haben. Ein digital verschicktes PDF ist für einen Menschen gut lesbar, für Software aber
oft nur ein Bild ohne strukturierten Inhalt. Handelt es sich gar um einen Scan oder ein Foto
einer Papierrechnung, ist der Informationsgehalt für eine Maschine erst recht nicht direkt
zugänglich.

Hinzu kommt die Heterogenität der Zielformate: Die PROGRAMMIERFABRIK GmbH benötigt Ausgaben in
ebInterface~6.1, dem österreichischen XML-Standard für den öffentlichen Bereich, und in
ZUGFeRD~2.3, dem deutschen Hybridformat, das ein PDF mit eingebettetem XML kombiniert.
Beide Formate unterliegen festen XML-Schemata (XSD), die für ein valides Dokument exakt
eingehalten werden müssen.

Die technischen Rahmenbedingungen sehen vor, dass das System als interne Webanwendung betrieben
wird. Ein Cloud-Deployment war nicht vorgesehen; der Betrieb läuft lokal über Docker.
Als KI-Backend kommt die externe Gemini~API (Google) zum Einsatz, die den rohen Rechnungstext
in strukturierte Daten überführt.

\section{Zielsetzung und Erfolgskriterien}
Ziel dieser Diplomarbeit war die Entwicklung eines funktionsfähigen Prototyps, der den
beschriebenen Verarbeitungsprozess vollständig abbildet. Konkret sollte das System in der
Lage sein:

\begin{enumerate}
    \item eine PDF-Rechnung hochzuladen und als valides ebInterface~6.1~XML auszugeben,
    \item ein Bild einer Rechnung per OCR zu verarbeiten und als valides ZUGFeRD~2.3~XML auszugeben,
    \item bereits verarbeitete Rechnungen persistent zu speichern und erneut herunterladen zu können.
\end{enumerate}

Als messbares Erfolgskriterium wurde definiert, dass die generierten XML-Dateien die offiziellen
Validatoren bestehen müssen: das ebInterface-Validierungsportal der AUSTRIAPRO\footnote{AUSTRIAPRO: \textit{ebInterface Validierungsportal}, \url{https://labs.ebinterface.at/}, letzter Zugriff am 20.01.2026}
und den deutschen E-Rechnungs-Validator für ZUGFeRD.\footnote{Vgl. Verein zur Förderung der Nutzung elektronischer Rechnungen e.V.: \textit{E-Rechnungs-Validator}, \url{https://erechnungs-validator.de/}, letzter Zugriff am 20.01.2026}

Dieses Kriterium wurde in den Abnahmetests erreicht: Zehn Testrechnungen, die die PROGRAMMIERFABRIK GmbH
zur Verfügung stellte, wurden vom System verarbeitet und lieferten XML-Ausgaben, die von beiden  verwebdeten Validatoren akzeptiert wurden.
Eine kurze Charakterisierung dieser zehn Testrechnungen (anonymisiert) findet sich im Anhang (Tabelle \ref{tab:testrechnungen}).

\section{Umfang, Annahmen und Abgrenzungen}
Der Umfang des Projekts umfasst folgende Komponenten:

\begin{itemize}
    \item \textbf{Backend}: RESTful Web API auf Basis von ASP.NET Core~9.0 mit
    PDF-Textextraktion (PdfPig), OCR (Tesseract), KI-Normalisierung (Gemini~API)
    und XML-Generierung für beide Zielformate.
    \item \textbf{Frontend}: Angular Single-Page-Application mit Upload-Komponente,
    Fortschrittsanzeige, Dokumentenvorschau und Download-Funktion.
    \item \textbf{Datenhaltung}: PostgreSQL-Datenbank mit Entity Framework Core (Code-First)
    zur persistenten Speicherung verarbeiteter Rechnungen.
\end{itemize}

Bewusst ausgeklammert wurden folgende Punkte:

\begin{itemize}
    \item \textbf{Produktionsreifes Cloud-Deployment}: Das System läuft lokal über Docker
    und ist nicht für den Betrieb auf einem öffentlich erreichbaren Server ausgelegt.
    \item \textbf{Weitere E-Rechnungsstandards} wie XRechnung oder Peppol~BIS,
    da diese außerhalb des unmittelbaren Bedarfs der PROGRAMMIERFABRIK GmbH liegen.
    \item \textbf{Benutzerverwaltung und Authentifizierung}, da es sich um ein rein
    internes Tool handelt.
    \item Weitere Features, die im Text explizit als \textit{nicht implementiert}
    gekennzeichnet sind.
\end{itemize}

Als Annahmen gelten, dass Eingangsrechnungen in deutscher oder englischer Sprache vorliegen
und die Gemini~API während des Betriebs erreichbar ist.

\section{Beiträge (Team- und Individualleistungen)}
Die Arbeit ist als Teamprojekt strukturiert. Gemeinsame Teile wie Einleitung,
Gesamtarchitektur, Technologiestack und Schlussbetrachtung wurden von beiden Autoren
gemeinsam erarbeitet. Die inhaltlichen Schwerpunkte sind dabei klar aufgeteilt.

\textbf{Sebastian Radić} verantwortet den Backend-Teil des Projekts.
Sein Beitrag umfasst die theoretischen Grundlagen zu E-Rechnungsstandards und Compliance
(ebInterface~6.1, ZUGFeRD~2.3, EN~16931), die vollständige Backend-Implementierung
(Controller, Services, Repository-Schicht, Datenbankzugriff) sowie die gesamte
Verarbeitungspipeline von der PDF-Extraktion über OCR und KI-Normalisierung
bis zur XML-Generierung.

\textbf{Luis Schörgendorfer} verantwortet den Frontend-Teil des Projekts.
Sein Beitrag umfasst die Anforderungsanalyse des Frontends, die Architektur- und
Technologieentscheidungen der Angular-Anwendung, die Implementierung der Upload-Komponente,
der UI/UX-Gestaltung und der Backend-Anbindung sowie die Optimierung der
Frontend-Performance und die abschließende Evaluation der Anwendung.

\section{Aufbau der Arbeit und Leseführung}
Die Arbeit ist in sieben Teile gegliedert, die den Entwicklungsprozess von der
theoretischen Grundlage bis zur abschließenden Bewertung nachvollziehbar machen.

\textbf{Teil~I -- Theoretische Grundlagen} erklärt die beiden E-Rechnungsstandards
ebInterface~6.1 und ZUGFeRD~2.3 sowie deren rechtlichen Rahmen.
Wer die Standards bereits kennt, kann diesen Teil überspringen.

\textbf{Teil~II -- Systemarchitektur und Design} beschreibt die Gesamtarchitektur
des \textit{SmartBillConverter}, den Datenfluss durch das System und die Auswahl
des Technologiestacks. Dieser Teil ist die Grundlage für alle folgenden
Implementierungskapitel.

\textbf{Teil~III -- Implementierung (Backend)} geht tief in den Code:
Backend-Struktur, Controller, PDF-Extraktion, OCR-Pipeline, KI-Normalisierung
und die Generierung beider XML-Formate werden Schritt für Schritt erklärt.

\textbf{Teil~IV -- Frontend-Anforderungsanalyse und Architektur} dokumentiert,
welche Anforderungen das Frontend erfüllen muss, wie das Datenmodell aussieht
und welche Technologieentscheidungen getroffen wurden.

\textbf{Teil~V -- Implementierung (Frontend \& Backend-Refinement)} zeigt die
konkrete Umsetzung der Angular-Komponenten, der UI/UX-Gestaltung und der
Anbindung ans Backend.

\textbf{Teil~VI -- Optimierung und Ergebnisse (Frontend)} bewertet die Performance
der Anwendung und fasst die Ergebnisse der Frontend-Evaluation zusammen.

\textbf{Teil~VII -- Ergebnisse, Diskussion und Ausblick} benennt die bekannten
Schwächen des Systems, diskutiert mögliche Erweiterungen und zieht ein
abschließendes Fazit.